\subsection{相对于子代数的谱}

在本小节中，我们假定 K = C。

引理 3. --- 设 \(X_{1}\) 与 \(X_{2}\) 是 C 的紧子集。若 \(X_{2}\) 包含在 \(X_{1}\) 中，且 \(X_{1}\) 在 C 中的边界包含在 \(X_{2}\) 中，则 \(X_{1}\) 是 \(X_{2}\) 与 \(X_{2}\) 在 C 中的补集的某些有界连通分支的并。

设 U 是 \(\mathbf{C}-X_{2}\) 的一个连通分支。\(X_{1} \cap U\) 在开集 U 中的每个边界点也是 \(X_{1}\) 在 C 中的边界点，故由假设属于 \(X_{2}\)；由于 \(U \cap X_{2} = \varnothing\)，我们看出 \(X_{1} \cap U\) 在空间 U 中没有边界点。由于 U 连通，交 \(X_{1} \cap U\) 或为空或等于 U (TG, I, p.~82, 推论)，由此得引理。

命题 6. — 设 A 是一个幺 Banach 代数，B 是 A 的一个闭幺子代数。对每个 \(x \in B\)，有 \(\mathrm{Sp}_{\mathrm{B}}(x) \supset \mathrm{Sp}_{\mathrm{A}}(x)\)，且 \(\mathrm{Sp}_{\mathrm{A}}(x)\) 在 \(\mathbf{C}\) 中的边界包含 \(\mathrm{Sp}_{\mathrm{B}}(x)\) 在 \(\mathbf{C}\) 中的边界。特别地，若 \(\mathrm{Sp}_{\mathrm{B}}(x) \subset \mathbf{R}\)，则我们有 \(\mathrm{Sp}_{\mathrm{B}}(x) = \mathrm{Sp}_{\mathrm{A}}(x)\)。

我们有 \(\mathrm{Sp}_{\mathrm{B}}(x) \supset \mathrm{Sp}_{\mathrm{A}}(x)\) (I, p.~3 的注 6)。若 \(\lambda\) 是 \(\mathrm{Sp}_{\mathrm{B}}(x)\) 在 C 中的边界上的一个点，则存在一个序列 \((\lambda_{n})\)（由 \(\mathrm{Sp}_{\mathrm{B}}(x)\) 的外部点组成），它趋于 \(\lambda\)。于是 \(x - \lambda_{n}\) 在 B 中可逆，且趋于 \(x - \lambda\)，后者在 B 中不可逆；因而 \(x - \lambda\) 是 B 中的一个左或右拓扑零因子 (I, p.~24 的命题 4)，从而也是 A 中的。于是 \(\lambda \in \mathrm{Sp}_{\mathrm{A}}(x)\)。但由于 \(\mathrm{Sp}_{\mathrm{A}}(x) \subset \mathrm{Sp}_{\mathrm{B}}(x)\)，复数 \(\lambda \in \mathrm{Fr}_{\mathbf{C}}(\mathrm{Sp}_{\mathrm{B}}(x))\) 不可能是 \(\mathrm{Sp}_{\mathrm{A}}(x)\) 的内点，故属于其边界。

推论. --- 集合 \(\mathrm{Sp}_{\mathrm{B}}(x)\) 是 \(\mathrm{Sp}_{\mathrm{A}}(x)\) 与 \(\mathbf{C}-\mathrm{Sp}_{\mathrm{A}}(x)\) 的某些有界连通分支的并。

这由命题 6 与引理 3 得出。

这个推论将由 I, p.~46 的命题 13 与 I, p.~46 的命题 14 加以补充。

